\chapter{Sprint 13 --- Mobile UI: Lịch sử đặt hàng (List + Detail) (2026-06-10)}

\section{Bối cảnh}

BA phát hành Figma FINAL cho trang \textbf{Lịch sử đặt hàng} mobile --- frame
\texttt{4445:4} (\texttt{Wujia\_Mobile\_Order\_History\_List\_Detail\_v10\_FINAL}),
file copy \texttt{aoeiDYlg6vlhJZg2w6Q7o5}. Frame là 2 ảnh phẳng đã export:
màn \textbf{List} (\texttt{/portal/purchase-history}) + màn \textbf{Detail}
(\texttt{/portal/purchase-history/<id>}).

Trước sprint, route \texttt{/portal/purchase-history} \textbf{chưa có bản mobile} ---
chỉ 1 layout desktop kiểu bảng 7 cột, tràn/khó đọc \(<992\)px. Mục tiêu: dựng bản
mobile theo Figma, đồng bộ shell Sprint 12 (header cyan + store strip + bottom-nav),
desktop GIỮ NGUYÊN 100%.

\textbf{Chốt từ user (2026-06-10):}
\begin{itemize}
  \item \textbf{Status badge = UI tạm}: chỉ map nhãn ở mức UI cho mobile (nhãn giàu
        hơn của Figma: "Đã xử lý / Đang giao / Hoàn tất"). KHÔNG build logic suy ra
        "Đang giao" từ batch/delivery state --- để sprint sau.
  \item \textbf{Mobile-only}: desktop list+detail giữ nguyên cả layout lẫn nhãn
        \texttt{STATE\_BADGES} cũ.
  \item \textbf{Ghi chú khi đặt hàng}: luôn hiện card (khớp layout Figma); rỗng thì
        placeholder "Không có ghi chú." Dùng \texttt{portal\_note} (KHÔNG fallback
        \texttt{order.note} vì đó là Terms\&Conditions của Odoo).
\end{itemize}

\section{Field mapping (KHÔNG cần field backend mới)}

\begin{longtable}{@{}lp{8.6cm}@{}}
\toprule
\textbf{UI} & \textbf{Field / nguồn} \\
\midrule
\endhead
Mã đơn & \texttt{order.name} (trigram index \(\rightarrow\) search \texttt{ilike} nhanh) \\
Giờ \(\cdot\) ngày & \texttt{date\_order.strftime('\%H:\%M · \%d/\%m/\%Y')} \\
Tổng tiền / Thành tiền & \texttt{amount\_total} / \texttt{line.price\_subtotal} \\
Người đặt & \texttt{portal\_requester\_user\_id.name or user\_id.name} \\
Cửa hàng & \texttt{franchise\_id.code} + " - " + \texttt{.name} \\
Số dòng hàng & \texttt{'\%02d' \% len(order.order\_line)} \\
Số phiếu batch & \texttt{batch\_id.name} (Many2one \texttt{stock.picking.batch}, store+index) \\
Thời gian khởi hành & \texttt{batch\_id.planned\_departure} (Datetime, \texttt{wujia\_delivery}) \\
Ghi chú & \texttt{portal\_note} (rỗng \(\rightarrow\) placeholder) \\
Quy cách & \texttt{line.product\_id.uom\_id.name} (không có field packaging riêng) \\
SL & \texttt{line.product\_uom\_qty} \\
\bottomrule
\end{longtable}

\section{Module \& file chạm}

\begin{itemize}
  \item \texttt{wujia\_portal\_purchase\_history}:
    \begin{itemize}
      \item \texttt{controllers/portal.py}: thêm \texttt{MOBILE\_STATE\_BADGES}
            (label map UI-only, tách khỏi \texttt{STATE\_BADGES} desktop) + param
            \texttt{q} (search Mã đơn, \texttt{ilike} trên \texttt{name}) + truyền
            \texttt{mobile\_state\_badges}/\texttt{q} vào context list + detail.
      \item \texttt{views/portal\_history.xml}: bọc toàn bộ markup desktop trong
            \texttt{d-none d-lg-block} (giữ nguyên) + thêm sibling
            \texttt{d-lg-none} cho cả List (\texttt{.wujia-mhist}) và Detail
            (\texttt{.wujia-mhist-detail}).
    \end{itemize}
  \item \texttt{wujia\_portal\_layout}:
    \begin{itemize}
      \item \texttt{\_variables.css}: cụm token \texttt{-{}-wujia-mhist-chip-*}
            (chip lọc); phần còn lại reuse palette \texttt{-{}-wujia-morder-*}.
      \item \texttt{\_components.css}: class \texttt{.wujia-mhist-*} (title /
            search / chip / listhead / list-row / pager / summary / card / kv /
            product-row / empty). Mobile-only, reuse \texttt{.wujia-badge-*}.
      \item \texttt{views/assets.xml}: bump \texttt{?v} 1116/1117
            \(\rightarrow\) \textbf{1120} rồi \textbf{1121} (sau fix kv align).
    \end{itemize}
\end{itemize}

\section{Status mapping (UI tạm)}

\begin{lstlisting}[language=Python]
# Mobile-only, TÁCH khỏi STATE_BADGES desktop. "shipping" để sẵn nhưng CHƯA
# có state sale.order nào trỏ tới -> wire batch state ở sprint sau.
MOBILE_STATE_BADGES = {
    'draft':    ('Mới',       'wujia-badge-muted'),
    'sent':     ('Đã gửi',    'wujia-badge-info'),
    'sale':     ('Đã xử lý',  'wujia-badge-success'),
    'shipping': ('Đang giao', 'wujia-badge-warning'),
    'done':     ('Hoàn tất',  'wujia-badge-success'),
    'cancel':   ('Đã hủy',    'wujia-badge-danger'),
}
\end{lstlisting}

\section{Gì đã làm}

Cửa hàng nhượng quyền nay xem \textbf{lịch sử đặt hàng} trên điện thoại gọn gàng
đúng thiết kế: màn danh sách có ô tìm theo mã đơn, lọc khoảng ngày + chip nhanh
(Tháng này / Tháng trước / Đã xử lý / Xóa lọc), mỗi đơn 1 dòng (mã \(\cdot\) giờ
ngày \(\cdot\) trạng thái \(\cdot\) tổng tiền). Bấm vào đơn mở màn chi tiết gồm 4
thẻ: tóm tắt, thông tin đơn hàng, thông tin batch/giao hàng, ghi chú, và danh sách
sản phẩm (quy cách + số lượng). Bản desktop không đổi.

\section{Lý do / Trade-off}

\begin{itemize}
  \item \textbf{Status UI tạm}: Figma có nhãn "Đang giao" --- không có state native
        nào của \texttt{sale.order} tương ứng (cần ghép thêm trạng thái
        batch/picking). Theo chốt của user, sprint này chỉ map nhãn UI; logic suy
        ra "Đang giao" để sau \(\rightarrow\) mobile dùng \texttt{MOBILE\_STATE\_BADGES}
        riêng, desktop giữ \texttt{STATE\_BADGES} cũ.
  \item \textbf{Search "Mã" wire luôn}: tuy là sprint UI, vẫn nối \texttt{q} (1 dòng
        domain \texttt{ilike}) để ô tìm không "chết"; an toàn perf nhờ trigram index
        sẵn có trên \texttt{sale.order.name}.
  \item \textbf{Gutter}: \texttt{.wujia-mhist} không tự pad ngang --- dựa
        \texttt{.content-wrapper} (Vuexy 2.2rem) như \texttt{.wujia-morder} để khớp
        gutter trang Đặt hàng.
  \item \textbf{kv value căn trái}: ban đầu căn phải sát lề; user feedback
        \(\rightarrow\) đổi sang cột cố định (label \texttt{flex 0 0 110px}, value
        \texttt{text-align:left}) đúng Figma.
\end{itemize}

\section{Verify}

Upgrade RC=0; \texttt{test\_sprint9} 7/7 + \texttt{test\_sprint5} 20/20 (regression
sạch); HTTP authenticated render List + Detail = 200, 0 traceback, đủ marker
\texttt{wujia-mhist-*}; định dạng khớp Figma (\texttt{155.250 đ},
\texttt{16:50 · 27/05/2026}, \texttt{01 sản phẩm}); desktop markup (bảng +
state-step) còn nguyên; CSS served \texttt{?v=1121}. Visual \(<992\)px do user xác
nhận trên browser thật.
